\section{Introduction}\label{sec:intro}
The unit-cost edit distance $\ED(x,y)$ is the minimum number of
single-symbol insertions, deletions and substitutions that transform a string
$x$ into a string $y$. All three operations have cost one. The strings are
stored explicitly; write $N=|x|+|y|$. Their symbols are integer labels of
polynomial magnitude in $N+2$, stored in words of $O(\log(N+2))$ bits.
There is no bound on the alphabet size independent of $N$.

An approximation scheme receives a rational $\varepsilon\in(0,1)$ and
estimates the distance multiplicatively. Its guarantee must therefore include
returning zero when the two strings are equal. We prove the following result.

\begin{theorem}\label{thm:main}
There is one uniform randomized algorithm that, given $x,y$ and a rational
$\varepsilon\in(0,1)$, returns a nonnegative integer $\widehat D$ such that
\[
 \Pr\bigl[\ED(x,y)\le \widehat D\le(1+\varepsilon)\ED(x,y)\bigr]
 \ge \frac23.
\]
For every fixed rational $\varepsilon$, its worst-case expected running time
on strings of total length $N$ is $N^{1+o(1)}$ as $N\to\infty$ on the
logarithmic-word RAM. The expectation is over internal randomness only,
and the bound includes input processing, preprocessing, numerical work,
random sampling and all data-structure operations. Equal strings return zero
deterministically.
\end{theorem}

The program is uniform in the accuracy input. The asymptotic assertion is
pointwise in each fixed rational accuracy: its constants and the threshold
at which the bound applies may depend on $\varepsilon$. We do not assert
polynomial dependence on $1/\varepsilon$ or a joint bound when the accuracy
varies with $N$. The output is a distance estimate; no edit-script
reconstruction is needed.

The algorithm uses exact computation below an accuracy-dependent input-length
threshold. Section~\ref{sec:geometry} gives the finite input convention and
fallback rules; their thresholds are very large, so the result is an
asymptotic approximation scheme.

\subsection{Context and approach}
Edit distance originates in the study of errors in transmitted strings,
including Levenshtein's binary codes for insertions, deletions and bit
flips~\cite{Levenshtein1966}. The dynamic program of Wagner and
Fischer~\cite{WagnerFischer1974} computes the minimum cost of insertions,
deletions and substitutions in quadratic time. Masek and
Paterson~\cite{MasekPaterson1980} applied block tabulation to obtain a
logarithmic speedup for fixed alphabets and suitable fixed costs.
Distance-sensitive exact algorithms offer another route when few edits
are needed: Ukkonen~\cite{Ukkonen1985} developed bounded-distance methods,
and Landau and Vishkin~\cite{LandauVishkin1988,LandauVishkin1989} combined
furthest-reaching states with fast matching queries. Myers's related
difference algorithm~\cite{Myers1986} treats insertion--deletion scripts.
For unrestricted exact computation, Backurs and
Indyk~\cite{BackursIndyk2018} showed that a strongly subquadratic
deterministic algorithm
would refute the Strong Exponential Time Hypothesis. That lower bound
does not rule out fixed-factor approximation.

Approximation algorithms progressively improved the accuracy available
with substantially less work. In the following historical comparisons,
write $n$ for a common upper bound on the string lengths.
Batu and coauthors~\cite{BatuEtAl2003}
gave sublinear tests for fixed alphabets that distinguish sufficiently
small distance from linear distance. For general pairs, Bar-Yossef, Jayram, Krauthgamer and
Kumar~\cite{BarYossefEtAl2004} obtained an $n^{3/7}$ approximation in
quasi-linear time; Batu, Erg\"un and
Sahinalp~\cite{BatuErgunSahinalp2006} improved the factor to
$n^{1/3+o(1)}$. For binary strings, Ostrovsky and
Rabani~\cite{OstrovskyRabani2007}
constructed a low-distortion embedding into $\ell_1$, and Andoni and
Onak~\cite{AndoniOnak2012} obtained an efficiently computable
$2^{\widetilde O(\sqrt{\log n})}$ approximation in $n^{1+o(1)}$ time.
Andoni, Krauthgamer and Onak~\cite{AndoniKrauthgamerOnak2010} then
achieved a $(\log n)^{O(1/\xi)}$ factor in $n^{1+\xi}$ time for each
fixed $\xi>0$, using a hierarchical alignment representation and
nonuniformly allocated precisions.

The next advances reached constant factors. Chakraborty, Das,
Goldenberg, Kouck\'y and Saks~\cite{ChakrabortyEtAl2020} gave a
constant-factor approximation in $\widetilde O(n^{12/7})$ time.
Brakensiek and Rubinstein~\cite{BrakensiekRubinstein2020} and Kouck\'y
and Saks~\cite{KouckySaks2020} independently obtained near-linear
exponents for sufficiently far input pairs, with a material additive
error on general pairs. Andoni and Nosatzki~\cite{AndoniNosatzki2020}
removed that restriction: for every fixed $\xi>0$, their algorithm
gives an approximation factor depending only on $\xi$ in
$O(n^{1+\xi})$ time. The factor depends on the exponent parameter;
this is not an approximation scheme with error tending to zero.
For accuracy arbitrarily close to one, Mao and
Rubinstein~\cite[Theorem~6.1]{MaoRubinstein2026} recently obtained a
randomized $(1+\varepsilon)$ approximation in
$n^2/2^{\log^{\Omega(1)}n}$ time for every fixed positive
$\varepsilon$.

The issue addressed here is how to retain fixed $(1+\varepsilon)$
accuracy while reducing the expected total work to $N^{1+o(1)}$.
Recursive interval decompositions, economical evaluation of local
distances, and accounting along an optimal alignment have substantial
precedent in the algorithms above. In our refinement, unrestricted recursive evaluation of child distances would consume
the available time even if the resulting numerical error were small. We therefore need both a narrow family that
contains a suitable comparison alignment and a bound on the number of
child distances actually evaluated.

For the initial estimates we adapt the shift tree and precision sampling
of Andoni, Krauthgamer and
Onak~\cite{AndoniKrauthgamerOnak2010,AndoniKrauthgamerOnak2011}.
The refinement is related to Mao and Rubinstein's use of multiscale
deviation to constrain alignment paths before sampling
\cite{MaoRubinstein2026}. Here we measure the alignment's changes in relative position using
preliminary cost estimates as weights. A telescoping bound on their
variation controls the total cost of places where the prediction fails.
An entropy-regularized prediction rule, within the framework of
Rakhlin and Sridharan~\cite{RakhlinSridharan2013}, then controls the
accumulated error where the predictions can be used. We implement its
optimization by a finite Frank--Wolfe procedure
\cite{FrankWolfe1956,Jaggi2013}. The specialized geometric, probabilistic
and numerical guarantees are proved in this paper.

Two related comparisons concern different requirements on the output or
the edit costs. A single $\ell_1$ embedding of all binary strings up to a
length bound must incur growing distortion
\cite[Theorem~1.1]{OpenAIEditDistortion2026}. Our algorithm estimates the
distance of the supplied pair and does not construct such a common
embedding. Recent algorithms for metric-weighted edit costs obtain
$(3+\varepsilon)$ approximation in strongly subquadratic time
\cite{MaderTavasoliWang2026,DasKipouridisKociumaka2026}; their cost model
is broader than the unit-cost model treated here.

\subsection{Proof overview}
Recursively partition the first string into a tree of intervals, ending at
intervals of length at most one. An interval state consists of a source
interval $I$ and a target interval in the second string; its true cost is
the edit distance between the two substrings. An action at a nonleaf
assigns one target interval to each source child. We charge the sum of
the child estimates together with penalties for gaps or overlaps between
successive target endpoints and at the two outer endpoints. These
penalties let us compare nearby actions without requiring their target
intervals to join exactly.

\paragraph{Estimates that can be requested adaptively.}
The initial seed is an integer table whose entries give coarse upper
approximations to all interval-state costs. Although the algorithm evaluates
only requested entries, the accuracy event covers the entire domain:
later queries may depend on earlier random answers. Each initial entry
uses a private random stream, so its expected evaluation cost can also
be bounded conditional on the request that first exposes it. Accuracy
and expected work thus use different probability statements. Successive
refinement passes improve the seed approximation factor.

\paragraph{Predicting where an alignment goes.}
One refinement starts from such a seed. It first propagates estimates
through a wide family of actions and averages the resulting preliminary
tables to obtain cost proxies. At a source cut, the alignment's shift is
the target position minus the source position. The total shift change
divided by the predicted total cost suggests how this shift should evolve
as prefix cost accumulates. Requiring agreement with that prediction
restricts the possible target positions at each cut to a narrow band.
The algorithm computes these bands from the cost proxies alone.

To analyze them, fix optimal alignments and recursively round their child
endpoints. This gives a comparison tree that the algorithm never computes.
A cost-weighted variance measures how far the comparison shifts depart
from the predicted evolution. Its sum telescopes, with additional charges
for disagreements among the cost proxies. Consequently the total cost of
nodes with large discrepancies is controlled. At the other nodes, either
the substring-length difference already accounts for almost all predicted
cost, giving a direct error bound, or the computed band contains the
comparison action.

\paragraph{From a band to a better seed.}
Within a band, an entropy-regularized rule predicts an action using earlier
child values. Shared samples also compare rounded actions across nearby
representatives. The resulting estimates decrease with the refinement
index. After conditioning on complete child tables, the parent samples
retain their prescribed laws, and an online inequality bounds cumulative
excess over the comparison action.
Some intermediate estimates can be smaller than true cost. The analysis
therefore clips them from below at true cost, so an underestimate at one
node cannot cancel an error elsewhere. Summing
over the comparison tree gives a small expected positive error at any
fixed queried state, together with a deterministic multiplicative lower
bound. Independent repetitions, rescaling and medians turn these two
bounds into a simultaneous one-sided seed with a smaller approximation
factor. Repeating this improvement and making one final refinement gives
the desired answer.

\paragraph{Paying for recursive queries.}
A wide search or an optimization step may inspect many actions, but it
reuses previously gathered child values. Only shared group samples,
selected rounded actions and actual online draws require child values
at the current refinement index. Common rounded grids limit the number
of these calls to slightly more than the tree's branching factor.
Along a dependency path, a call either descends the tree at the same
index or decreases an index or pass. Counting these two kinds of steps
bounds the total expected work, including rational arithmetic and exact
sampling, by $N^{1+o(1)}$.

Section~\ref{sec:geometry} fixes the computational model, geometry and
parameters. Sections~\ref{sec:seed}--\ref{sec:bands} construct the tables.
Sections~\ref{sec:comparison} and~\ref{sec:regret} prove the geometric and
probabilistic estimates, and Section~\ref{sec:accuracy} assembles the
accuracy guarantee. Sections~\ref{sec:numerics} and~\ref{sec:work} implement
every operation and complete the running-time proof.
